% Draft for \subsubsection{Structural Characterisation and Binding-Mode Assumptions}
% Verify/add the provisional BibTeX keys before inclusion in the thesis.
% New provisional keys introduced in the audit bullets below:
% hardy_aptarank_2026, uniprot_consortium_uniprot_2025,
% sabrowski_use_2022, king_ndm1_2013, cheng_distinct_2023,
% gelinas_embracing_2016.

\paragraph{Reproducible target-selection record.}
The target was selected in a target-only development phase before candidate
scoring.  The following chronological record maps the selection process to the
seven tasks specified at the start of the study and records negative as well
as positive decisions.

\begin{itemize}
  \item \textbf{Review of the computational prioritisation and evaluation
  approach (task~1).}  We reviewed the AptaRank camera-ready workflow and the
  MRes evaluation plan before inspecting candidate scores
  \cite{hardy_aptarank_2026}.  This established that Tier~2 is a target-aware
  geometric-compatibility annotation rather than docking or binding
  prediction; that binding mode and site must be assigned from independent
  evidence; and that one structure--chain--mode--site configuration must be
  frozen before scoring.  The screen therefore did not use AptaRank score,
  fpocket rank, fpocket score or fpocket volume to select the target, and it did
  not run several configurations and retain the most favourable result.

  \item \textbf{Search and metadata extraction (task~2).}  On 9~October~2026,
  the RCSB PDB Search API v2 was queried for entries returned by the term
  ``NDM-1'', reproducing the 166-entry census visible in the RCSB interface
  \cite{berman_protein_2000}.  Entry, polymer-entity, resolution, mutation,
  title and non-polymer-component fields were then retrieved through the RCSB
  Data API (GraphQL), and coordinates were obtained from the RCSB file service.
  Polymer sequences were locally aligned to canonical \textit{Klebsiella
  pneumoniae} NDM-1, UniProtKB C7C422
  \cite{uniprot_consortium_uniprot_2025}.  The prespecified automated triage
  retained records with at least 98\% local sequence identity and 80\%
  canonical coverage (155/166), then identified 127 exact, possibly
  terminally truncated, canonical matches with no aligned mismatch.  For the
  native-like and carbapenem-product lists, these were further restricted to
  structures reported as unmutated and with resolution no poorer than
  2.5~\AA{}.  Title and heterogroup keyword classification was used only for
  triage; each shortlisted record and its primary publication was subsequently
  checked manually because an automated ``apo-like'' label can be confounded
  by crystallisation additives or incomplete ligand annotations.

  \item \textbf{Comparison and filtering categories.}  The surviving
  structures were compared in a fixed evidence table using: exact NDM-1
  identity and reported mutations; experimental method, resolution and
  $R_{\mathrm{free}}$; construct boundaries and author numbering; completeness
  of the mature catalytic domain and the mobile L3/ASL1 and L10/ASL4 regions;
  chain and biological-assembly assignment; native, substrate/product or
  inhibitor-bound state; identity and occupancy of catalytic metals; pH and
  crystallisation or soaking conditions; non-water heterogroups and their
  proximity to the catalytic groove; and resemblance to the soluble,
  recombinant NDM-1 state used in the reference aptamer selections.  Coordinate
  parsing checked residues 42--270, catalytic residues 120, 122, 124, 189,
  208, 211, 220 and 250, L3 residues 63--73 and L10 residues 210--228.
  C$\alpha$ conformations were compared after residue-number-matched Kabsch
  alignment, with whole-domain, L3 and L10 RMSDs reported separately.

  \item \textbf{Shortlisting.}  Five structures were retained for documented
  expert comparison, without scoring aptamers against any of them: PDB~5ZGX
  chain~A (wild-type, native-form, dizinc, pH~7.5, 0.95~\AA{}); PDB~3SPU
  (independent native dizinc comparator, 2.10~\AA{}); PDB~5XP6 chain~A
  (wild-type native structure at pH~5.5, 0.95~\AA{}); PDB~5YPL chain~A
  (hydrolysed-imipenem product complex, 1.80~\AA{}); and PDB~5N0H chain~A
  (hydrolysed-meropenem product complex, 1.90~\AA{})
  \cite{king_crystal_2011,king_ndm1_2013,zhang_active-site_2018}.  Product-bound
  5YPL and 5N0H were retained as mechanistic comparators rather than preferred
  targets because product binding stabilises a catalytic-groove conformation.
  PDB~3ZR9 was excluded from the primary set because of non-native metal
  substitution/additional metals; PDB~6C89 because it is a triple mutant;
  PDB~6TWT because L3 residues are unresolved; and PDB~4EYL because 5N0H is its
  later re-refinement.  Inhibitor-bound structures and redundant members of
  the 5ZGW/5ZGX/5ZGY/5ZGZ native series were not promoted to separate primary
  configurations.

  \item \textbf{Review of NDM-1 reference aptamers (task~3).}  The primary RNA
  evidence was von Laer's NDM-1 SELEX study, which selected against soluble,
  N-terminally His-tagged NDM-1 residues Gly29--Arg270 and identified D909 as
  the strongest of the reported RNA binders and an inhibitor in a nitrocefin
  assay \cite{ludwig_von_laer_selektion_2013}.  Sabrowski \textit{et al.}
  reported the ssDNA aptamer NDM1-H14.2\_43 selected against immobilised
  His--NDM-1 with His-tag masking and counterselection
  \cite{sabrowski_use_2022}.  Patents CN102965377B and CN108949766B were also
  reviewed for ssDNA sequences and claimed affinities; the latter deliberately
  selected a cross-reactive NDM-1/VIM-2 aptamer.  These ssDNA sources informed
  target-state plausibility but were not treated as direct evidence of an RNA
  epitope because RNA and ssDNA differ in sugar chemistry and accessible
  conformational ensembles.  None of the four sources supplied a
  residue-resolved NDM-1 aptamer contact map or competitive-binding structure.
  D909 inhibition and the thesis' proposed ``homing principle'' were therefore
  treated as suggestive, not as proof that D909 occupies the catalytic groove.

  \item \textbf{Review of NDM-1 structure, function and binding literature.}
  Structural and chemical-biology sources independently define a broad,
  solvent-exposed dizinc catalytic groove bordered by mobile loops.  The Zn1
  site is coordinated by His120, His122 and His189, and the Zn2 site by Asp124,
  Cys208 and His250; substrate/product recognition additionally involves the
  L3/ASL1 and L10/ASL4 regions
  \cite{kim_structure_2011,king_ndm1_2013,zhang_active-site_2018}.  The native
  active site is conformationally heterogeneous whereas product complexes are
  more consistent, and NMR work distinguishes apo-, dizinc- and
  inhibitor-associated states \cite{zhang_active-site_2018,cheng_distinct_2023}.
  General protein--aptamer structures show that RNA can recognise recessed
  grooves or extended surface patches \cite{gelinas_embracing_2016}; this did
  not override the absence of an experimentally mapped NDM-1 RNA epitope.

  \item \textbf{Final PDB ID and chain selection.}  PDB~5ZGX chain~A was fixed
  as the primary target because it is an unmutated, monomeric, native-form
  dizinc NDM-1 structure at near-physiological pH, determined at 0.95~\AA{}
  resolution, with residues 42--270 and both mobile catalytic-loop regions
  modelled \cite{zhang_active-site_2018}.  It is PDB entry version~1.3 and has
  the persistent identifier \texttt{10.2210/pdb5ZGX/pdb}.  Its bound ethylene
  glycol is a crystallisation additive, not a biological substrate: this and
  ordinary waters were removed during preparation.  The structure is a static
  soluble-domain crystal snapshot and does not reproduce NDM-1's native
  membrane anchoring, periplasmic environment or full conformational ensemble.

  \item \textbf{Binding-site and binding-mode selection (task~4).}  Because no
  independent residue-resolved aptamer surface epitope was available, surface
  mode would have required an arbitrary patch.  Pocket mode was therefore
  assigned a priori, with ``pocket'' operationally meaning the shallow catalytic
  groove rather than a deeply enclosed small-molecule pocket.  The six metal
  ligands His120, His122, Asp124, His189, Cys208 and His250 formed the primary
  fpocket-selection anchor set.  Nine independently defined extended-groove
  residues (Leu65, Met67, Phe70, Val73, Trp93, Gln123, Lys211, Gly219 and
  Asn220) were fixed solely as a tie-break set.  A cavity had to overlap at
  least one residue from each zinc subsite; an unresolved tie failed closed.
  Thus, although up to five plausible residue sets had been allowed during
  method development, only one prespecified rule and one resulting site were
  carried forward, avoiding a best-of-multiple search.

  \item \textbf{Configuration, validation and documentation (tasks~5--6).}
  A deterministic preparation script generated a biological master and a
  protein-only fpocket/AptaRank input, executed fpocket, inventoried every
  cavity and applied the fixed overlap rule \cite{le_guilloux_fpocket_2009}.
  Cavity~1 was uniquely selected; its lining residues were Leu65, Met67, Val73,
  Trp93, His122, Gln123, Asp124, His189, Lys211, Ala215, Ser217, Gly219,
  Asn220, His250 and Ser251.  The configuration, preparation policy, complete
  fpocket output, run report and SHA-256 manifest were frozen in
  \texttt{target/artifacts/ndm1\_5zgx\_fpocket4/}.  Seven unit tests passed, all
  manifest entries verified, and an independent replication recovered the same
  pocket, lining residues, 42 alpha spheres and 9.212~\AA{} characteristic
  dimension.  The Monte Carlo volume differed by 5.0\% and was retained only as
  a descriptive check.  No aptamer candidate, candidate score or AptaRank run
  was read or executed during this target-selection phase.

  \item \textbf{Time-box and deferred scope (task~7).}  The one-hour request
  was used as a prioritisation constraint: public metadata were screened
  programmatically, detailed review was limited to a five-structure shortlist,
  and one defensible configuration was frozen.  The time-box was not used to
  suppress failed checks or to assert unverified elapsed time.  Alternative
  structures, surface epitopes, allosteric sites, metal/water treatments and
  multi-configuration sensitivity testing are explicitly deferred to the PhD;
  they must not be searched post hoc for a favourable candidate result.

  \item \textbf{Software, services and version record.}  Scientifically
  material tools called during this work were: RCSB Search API v2, RCSB Data
  API/GraphQL and RCSB coordinate service (service build not exposed; accessed
  9~October~2026); UniProtKB REST for C7C422 (service build not exposed;
  accessed 9~October~2026); Python~3.13.3 and NumPy~2.4.6 for exploratory
  metadata, coordinate and Kabsch analyses; and, for the frozen workflow,
  Python~3.10.18, NumPy~1.26.4 and fpocket~4.0 in a conda~25.7.0 environment.
  Source PDFs were text-extracted with pypdf~6.10.0 and visually checked after
  rendering with Poppler \texttt{pdfinfo}/\texttt{pdftoppm}~26.05.0.  Files
  were retrieved with curl~8.7.1; SHA-256 values used
  \texttt{shasum}~6.02; repository inspection used Git~2.39.5 and
  ripgrep~15.2.0.  Python's standard-library \texttt{unittest} and
  \texttt{py\_compile} under Python~3.10.18 supplied validation and syntax
  checks.  Homebrew~7.0.3 was queried only for package availability and did not
  contribute analytical output.  The AptaRank implementation reviewed was
  version~0.1.0 (target schema~1.0.0), with ViennaRNA~2.7.2,
  forgi~2.2.3 and fpocket~4.0; AptaRank, ViennaRNA, forgi, freeSASA and APBS
  were not executed in this target-selection phase.  The built-in Codex web
  retrieval interface assisted source discovery and drafting; its service
  build was not exposed, so all scientific claims were checked against the
  cited primary sources or archived local documents.  Routine file-display
  utilities do not alter scientific output and were not separately versioned.
\end{itemize}

% The thesis preamble requires \usepackage{booktabs,tabularx} for this table.
\begin{table}[htbp]
  \centering
  \small
  \caption[NDM-1 structure candidates considered during target selection]{%
  Score-blind comparison of the five final NDM-1 structure candidates and
  ten near-shortlist structures retained in the screening audit. ``Complete''
  denotes that the catalytic residues and the L3/ASL1 and L10/ASL4 regions
  used in the coordinate check were modelled; it does not imply that the
  structure represents the full native conformational ensemble. Ligand codes
  are wwPDB chemical-component identifiers. Resolution and entry annotations
  were retrieved from RCSB PDB on 9~October~2026
  \cite{berman_protein_2000}.}
  \label{tab:ndm1-structure-screen}
  \begin{tabularx}{\textwidth}{@{}llclXX@{}}
    \toprule
    PDB ID / chain & Structural state & Resolution
      (\AA{}) & Catalytic region & Notable metals or heterogroups
      & Screening disposition \\
    \midrule
    \textbf{5ZGX / A} & Native, pH~7.5 & 0.95 & Complete
      & Two Zn, bridging OH; EDO in coordinates
      & \textbf{Selected primary target}; unmutated monomer, one chain and
        near-physiological pH \\
    3SPU / A--E & Apo/native comparator & 2.10 & Complete in five copies
      & Two Zn per chain; no non-water ligand other than Zn
      & Finalist; independent native comparator, but lower resolution and
        multiple crystallographic copies \\
    5XP6 / A & Native, pH~5.5 & 0.95 & Complete
      & Two Zn, OH, succinate and EDO; succinate contacts the groove
      & Finalist; high resolution, but acidic pH and active-site additive \\
    5YPL / A & Hydrolysed-imipenem EP product complex & 1.80 & Complete
      & Two Zn and hydrolysed imipenem (HIW)
      & Finalist mechanistic comparator; product-stabilised conformation \\
    5N0H / A & Hydrolysed-meropenem product complex & 1.90 & Complete
      & Two Zn and hydrolysed meropenem (8YL)
      & Finalist mechanistic comparator; preferred re-refinement of 4EYL \\
    \midrule
    5ZGW / A & Native, pH~7.5 after pH~5.5 soak & 0.95 & Complete
      & Two Zn, OH, EDO and glycerol
      & Near-shortlist; redundant native-series conformation with soaking
        history and additives \\
    5ZGY / A & Native, pH~7.5, Bis--Tris & 0.95 & Complete
      & Two Zn, OH and EDO
      & Near-shortlist; redundant member of the 5ZG series \\
    5ZGZ / A & Native, pH~7.5, imidazole & 0.95 & Complete
      & Two Zn, OH and EDO
      & Near-shortlist; redundant series member and imidazole condition \\
    5ZGI / A & Native, pH~6.5, succinate & 0.98 & Complete
      & Two Zn, OH, succinate and EDO; succinate contacts the groove
      & Near-shortlist; lower pH and active-site additive \\
    5ZR8 / A, E & Native-like & 1.50 & Complete
      & Chain A contains four Zn; chain E contains two Zn; sulfate present
      & Excluded from final five because metal/chain choice was less
        biologically direct \\
    5YPN / A & Hydrolysed-meropenem EI2 intermediate & 2.12 & Complete
      & Two Zn and meropenem intermediate (LMP)
      & Mechanistic comparator only; ligand-selected intermediate state \\
    4EYL / A & Hydrolysed-meropenem product complex & 1.90 & Complete
      & Two Zn and hydrolysed meropenem (0RV)
      & Excluded in favour of its later re-refinement, 5N0H \\
    6TWT / A, B & Nominally native, N-terminally truncated & 0.95
      & L3 incomplete: A lacks 66--71; B lacks 68--70
      & Two Zn; calcium and chloride present
      & Excluded because the pocket-bordering L3 loop is unresolved \\
    3ZR9 / A & Native-like metal-bound structure & 1.91 & Complete
      & One Zn plus Cd, Co and Ni; Cd occupies the second metal region
      & Excluded because the catalytic metal environment is non-native \\
    6C89 / A & Mutant mechanistic structure & 1.75 & Modelled; three mutations
      & Zn and sodium reported
      & Excluded because a triple mutant was not an appropriate wild-type
        prioritisation target \\
    \bottomrule
  \end{tabularx}
\end{table}

% Stable source identifiers for provisional bibliography entries:
% Sabrowski et al. (2022): 10.1038/s41598-022-12062-2.
% King et al. (2013): 10.1371/journal.pone.0047249; PMID 23363572.
% Zhang et al. (2018): 10.1128/AAC.01579-18; PMID 30150473.
% Cheng et al. (2023): 10.1021/jacsau.2c00651; PMID 37006760.
% Kim et al. (2011): 10.1371/journal.pone.0024621.
% Le Guilloux et al. (2009): 10.1186/1471-2105-10-168.

The target configuration was fixed before candidate scoring and without access
to candidate-specific AptaRank results. We selected the wild-type, native-form
NDM-1 structure PDB~5ZGX, chain~A, determined by X-ray crystallography at
0.95~\AA{} resolution at pH~7.5 \cite{berman_protein_2000,zhang_active-site_2018}.
The deposited model contains residues 42--270 in PDB author numbering and two
catalytic zinc ions. Alternate atom locations were resolved deterministically
by selecting the highest-occupancy conformer, with ties preferring a blank
alternate-location identifier and then conformer~A. A biological preparation
master retained the protein, the two zinc ions and the deposited bridging
hydroxide, while ethylene glycol and ordinary crystallographic waters were
removed. The geometry input used by AptaRank contained chain-A protein atoms
only, corresponding to the target configuration option
\texttt{strip\_hetatm: true}. This prevents active-site ethylene glycol from
being treated as part of NDM-1 while retaining the deposited metal-bound
protein conformation.

NDM-1 was assigned to AptaRank's pocket-binding mode, with ``pocket''
operationally encompassing its shallow, solvent-exposed catalytic groove. The
available NDM-1 RNA-aptamer studies do not map an aptamer epitope at residue
resolution; the active-site assignment is therefore an explicit modelling
hypothesis rather than a claim that the reference aptamer is known to bind
competitively at this site. Pocket mode was preferred over surface mode because
the binuclear catalytic groove is independently defined by structural and
functional evidence, whereas no independently mapped NDM-1 aptamer surface
patch was available.

fpocket version~4.0 was run once on the prepared protein-only structure using
300 Monte Carlo volume iterations \cite{le_guilloux_fpocket_2009}. Before the
run, the pocket-selection anchors were fixed as the six metal-coordinating
residues His120, His122, Asp124, His189, Cys208 and His250. The selected cavity
was the fpocket cavity with the greatest number of overlaps with this set. A
predefined extended-groove set (Leu65, Met67, Phe70, Val73, Trp93, Gln123,
Lys211, Gly219 and Asn220) was reserved only to resolve a tie. A valid cavity
was additionally required to contain at least one residue associated with each
zinc subsite; an unresolved tie or failure to cover both subsites would have
invalidated the configuration rather than triggering selection of fpocket's
highest-scoring unrelated cavity.

fpocket identified ten cavities. The prespecified rule uniquely selected
cavity~1, whose lining-residue set was Leu65, Met67, Val73, Trp93, His122,
Gln123, Asp124, His189, Lys211, Ala215, Ser217, Gly219, Asn220, His250 and
Ser251. It overlapped four of the six primary anchors, including residues from
both zinc subsites. Following the AptaRank workflow, the characteristic pocket
dimension was calculated from the largest principal-axis extent of the 42
fpocket alpha-sphere centres and was 9.212~\AA{}. The other principal-axis
extents were 2.721 and 0.890~\AA{}. The frozen run's fpocket volume estimate was
590.122~\AA$^3$, corresponding to an equivalent-sphere diameter of
10.407~\AA{}; the ratio between this diameter and the primary extent was 1.130,
so the predefined odd-shape warning was not raised. Because fpocket estimates
volume by Monte Carlo sampling, volume was treated only as a descriptive
cross-check and did not enter pocket selection. An independent replication
recovered the same cavity, lining residues, alpha-sphere count and
characteristic dimension exactly; its volume differed by 5.0\%, consistent
with the stochastic estimator.

The resulting structure--chain--mode--site configuration and its preparation
policy, software versions, raw and derived coordinate hashes, complete fpocket
output and pocket inventory were frozen before aptamer scoring. This
configuration measures compatibility between an estimated accessible RNA-loop
extent and one static representation of the NDM-1 catalytic groove. It does
not constitute molecular docking and does not establish binding, affinity,
competitive inhibition, selectivity, an RNA contact map or a complete
RNA--protein pose. Alternative NDM-1 conformations, surface epitopes,
allosteric sites and binding modes were not selected according to candidate
scores and are deferred to future work.
